%% kapa: exam=1 round=37 year=2026 period=1 form=A subject=real_estate subject_ko=부동산학원론 seq=14 number=94
%% source: Q-Net 감정평가사 기출문제 (cst00309) artlSeq=5258924 file=raw/1008/5258924/2026년 제37회 감정평가사 1차 1교시.pdf
%% conversion: kordoc 4.24.51 -> tools/gen1.py (auto); review=proofread-sonnet; flags=none
\begin{question}{94}
도시와 도시화에 관한 설명으로 옳지 않은 것은?
\begin{choices}
  \item 일반적으로 도시화의 4단계는 도시화(urbanization), 교외화(suburbanization), 탈도시화(disurbanization), 재도시화(reurbanization)로 설명된다.
  \item 도시화를 측정하는 지표로는 인구지표, 토지이용지표, 농촌요소의 쇠퇴와 도시요소의 증대를 나타내는 경제활동지표 등을 활용한다.
  \item 도시스프롤(urban sprawl) 현상은 도시의 팽창에 따라 도시가 무질서하게 교외로 확대되어 가는 현상을 말한다.
  \item 가도시화(pseudo-urbanization)는 도시의 인구가 도시 밖으로 벗어나 감소하는 현상을 말한다.
  \item 도시재생(urban regeneration)은 낙후된 기존 도시지역에 새로운 기능을 도입하거나 창출함으로써 경제적ㆍ사회적ㆍ물리적 환경을 향상시키는 활동을 말한다.
\end{choices}
\end{question}
\answer{4}
